% Autogenerated translation of LANGUAGE_PLAN.md by Texpad
% To stop this file being overwritten during the typeset process, please move or remove this header

\documentclass[12pt]{book}
\usepackage{graphicx}
\usepackage{fontspec}
\usepackage[utf8]{inputenc}
\usepackage[a4paper,left=.5in,right=.5in,top=.3in,bottom=0.3in]{geometry}
\setlength\parindent{0pt}
\setlength{\parskip}{\baselineskip}
\setmainfont{Helvetica Neue}
\usepackage{hyperref}
\pagestyle{plain}
\begin{document}

\chapter*{Contract language — v1 Roadmap and Implementation Plan}

This document outlines concrete plans to evolve the Contract language tooling at three different levels. It gives options, estimated effort, deliverables, milestones, risks, and recommended next steps so we can pick a pragmatic path forward.

\section*{Goals}

\begin{itemize}
\item Ship stable v1 language semantics and tooling useful for editor integration, testing, and the reference VM.
\item Provide precise diagnostics and editor features (diagnostic spans, hover signatures, completions) quickly.
\item Maintain a single canonical implementation where practical (the C compiler/runtime) while allowing pragmatic shortcuts for editor tooling.
\end{itemize}

\hrule
\section*{Options overview}

A — Near-term: Grammar + parser generator (recommended for editor UX)

\begin{itemize}
\item What: Convert the EBNF in \texttt{LANGUAGE\_SPEC\_v1.md} into a PEG or nearley grammar and generate a JS parser used by the LSP server. The parser will produce ASTs and diagnostics with precise start/end spans.
\item Why: Fastest path to high-quality editor diagnostics and features. Keeps the C compiler as a runtime/IL emitter and uses the JS parser for LSP workflows.
\item Deliverables:

\begin{itemize}
\item \texttt{lsp/grammar/contract.ne} (nearley grammar) + generated parser files
\item \texttt{lsp/parser.js} wrapper that returns JSON AST + diagnostics for a given document
\item Updated \texttt{lsp/server.js} to prefer \texttt{lsp/parser.js} for diagnostics \& symbols
\item Tests under \texttt{tests/} verifying parser diagnostics and AST shape
\end{itemize}
\item Timeline: \textasciitilde{}1–2 days (prototype) — adjust depending on grammar edge cases.
\item Risks:

\begin{itemize}
\item Duplicate parsing logic (JS + C) requires periodic sync.
\item Adds Node dependency for parsing; acceptable because LSP already runs in Node.
\end{itemize}
\end{itemize}

B — Medium-term: Improve the native C parser (canonical path)

\begin{itemize}
\item What: Expand \texttt{src/compiler.c} to a robust recursive-descent parser with full precedence, proper AST node ranges, block scoping, and clearer symbol tables. The compiler will emit debug info mapping bytecode back to source spans.
\item Why: Single, canonical parser for compiler, IL emitter, and tooling that avoids duplication.
\item Deliverables:

\begin{itemize}
\item \texttt{src/parser.h} / \texttt{src/parser.c} refactor with parser modules and AST with start/end positions
\item Update \texttt{diagnose\_source\_str} to use AST ranges and richer diagnostics
\item Emission of debug/source-mapping into the bytecode (optional) and tests
\end{itemize}
\item Timeline: \textasciitilde{}3–7 days depending on desired completeness.
\item Risks:

\begin{itemize}
\item More C-level engineering; slower iteration and potential for regressions.
\item Larger changes to the build and test cycles.
\end{itemize}
\end{itemize}

C — Test-first: Conformance test-suite and CI (stabilize semantics)

\begin{itemize}
\item What: Build a comprehensive test-suite of positive/negative \texttt{.ct} examples and expected diagnostics/IL results, plus scripts and a lightweight comparison harness.
\item Why: Provides a safety net for evolving the parser/VM and is low-risk to implement.
\item Deliverables:

\begin{itemize}
\item \texttt{tests/} expanded with \texttt{good/} and \texttt{bad/} cases and expected \texttt{.diag.json} or \texttt{.il} outputs
\item \texttt{tests/run\_tests.sh} / \texttt{ci/check.sh} that run the suite and assert results
\item Optionally wire a CI job (GitHub Actions) that runs the test suite on each PR
\end{itemize}
\item Timeline: \textasciitilde{}1–2 days for a good starter suite.
\item Risks: none significant; this is a strong enabler for either A or B.
\end{itemize}

\hrule
\section*{Recommendation}

If your priority is editor UX (fast precise diagnostics + hover/completions), choose A (nearley/PEG parser in the LSP). It gives the best value quickly and integrates well with the existing Node LSP server. If you instead want a single canonical implementation in C and have time, choose B. In either case, adding C (conformance tests) is highly recommended before major refactors.

My suggested execution plan if you pick A now:

\begin{enumerate}
\item Convert EBNF -$>$ nearley grammar and implement \texttt{lsp/parser.js}.
\item Wire \texttt{lsp/server.js} to call the parser for diagnostics/symbols (fall back to \texttt{build/Contract} if the parser throws or for IL generation).
\item Add tests in \texttt{tests/} that assert parser diagnostics for known errors (e.g., missing paren, unmatched braces, type in header, etc.).
\item Iterate until diagnostics are stable and editor highlights are precise.
\end{enumerate}

Milestones \& Targets (A)
- M1 (day 1 morning): Nearley grammar draft; parse simple files (functions, contracts, calls).
- M2 (day 1 afternoon): Diagnostic generation with precise spans; integrate into LSP diagnostics publishing.
- M3 (day 2): Symbol extraction (function signatures), completions, hover support; tests green.

Success criteria
- LSP publishes precise start/end diagnostic ranges for all samples in \texttt{tests/}.
- Hover shows function signatures for functions in the document and imported files.
- The C compiler remains the canonical emitter and is not required for editor diagnostics (but still used for build/run).

\hrule
\section*{Next steps}

Reply with your choice: A, B or C (or a combination like A+C). Once you pick, I'll:

\begin{itemize}
\item expand the todo list with fine-grained tasks,
\item implement the first milestone and run the test-suite,
\item open a change with the updated artifacts and show how to run them locally.
\end{itemize}

If you want A (nearley), I’ll start converting the EBNF to a nearley grammar immediately.

\end{document}
